\begin{figure}[htbp]
\centering
\begin{tikzpicture}[x=1cm,y=1cm,font=\small,
  >={Stealth[length=1.7mm]},line cap=round,
  nucleus/.style={circle,draw=figInk,fill=white,line width=.55pt,
    minimum size=5mm,inner sep=0pt,font=\small},
  factor/.style={draw=figInk!50,rounded corners=3pt,
    minimum width=1.05cm,minimum height=.57cm,inner sep=3pt}]
\path[use as bounding box] (0,-.35) rectangle (14,5.55);
\fill[figInk!3,rounded corners=7pt] (.4,2.55) rectangle (13.6,5.45);
\node[anchor=west,font=\scriptsize,text=figMuted] at (.67,5.13) {configuration columns};
\foreach \xx/\leftlabel/\rightlabel/\leftfill/\rightfill in
  {3.1/1/2/figHydrogen/figHydrogen,10.9/2/1/figHydrogen/figHydrogen} {
  \begin{scope}[shift={(\xx,4.08)}]
    \draw[figInk,line width=.75pt] (-.85,-.15)--(0,.62)--(.85,-.15);
    \node[nucleus,fill=\leftfill] at (-.85,-.15) {$\mathrm H_{\leftlabel}$};
    \node[nucleus,fill=\rightfill] at (.85,-.15) {$\mathrm H_{\rightlabel}$};
    \node[nucleus,fill=figOxygen,text=white] at (0,.62) {$\mathrm O_3$};
  \end{scope}
}
\draw[figInk,->,line width=.8pt] (5.25,4.15)--(8.75,4.15)
  node[midway,above,inner sep=6pt] {$\sigma=(12)$};
\node at (3.1,2.98) {$X=(\mathbf x_1,\mathbf x_2,\mathbf x_3)$};
\node at (10.9,2.98) {$\sigma X=(\mathbf x_2,\mathbf x_1,\mathbf x_3)$};
\node at (3.10,2.13) {$\Psi(X)$};
\node at (10.90,2.13) {$\Psi(\sigma X)$};
\node[factor,fill=figHydrogen] at (1.78,1.47) {$\xi_1$};
\node at (2.44,1.47) {$\otimes$};
\node[factor,fill=figHydrogen] at (3.10,1.47) {$\xi_2$};
\node at (3.76,1.47) {$\otimes$};
\node[factor,fill=figOxygen,text=white] at (4.42,1.47) {$\xi_3$};
\node[factor,fill=figHydrogen] at (9.58,1.47) {$\xi_2$};
\node at (10.24,1.47) {$\otimes$};
\node[factor,fill=figHydrogen] at (10.90,1.47) {$\xi_1$};
\node at (11.56,1.47) {$\otimes$};
\node[factor,fill=figOxygen,text=white] at (12.22,1.47) {$\xi_3$};
\node[draw=figGold,fill=white,text=figGold,rounded corners=3pt,
  minimum height=.57cm,inner sep=4pt] at (8.20,1.47) {$\pm$};
\node at (8.92,1.47) {$\bigl($};
\node at (12.90,1.47) {$\bigr)$};
\draw[figInk,->,line width=.8pt] (5.25,1.47)--(7.45,1.47)
  node[midway,above,inner sep=6pt] {$\stat(\sigma)\,\sigma$};
\node at (7,.28) {$\stat(\sigma)=(-1)^{2I}
  =\begin{cases}
    +1 & \text{integer spin number},\\
    -1 & \text{half-integer spin number}.
  \end{cases}$};
\end{tikzpicture}
\caption{For $\sigma=(12)$, the physical-state condition relates
the spin vectors at $X$ and $\sigma X$ by exchanging factors and
multiplying the whole vector by $\stat(\sigma)$. Here $I$ is the common
spin number of nuclei $1$ and $2$. A product value is shown, and the
rule extends linearly to all spin vectors. The relation holds almost
everywhere, as expressed by $\Psi\circ\sigma\sim\stat(\sigma)\,\sigma\act\Psi$.}
\label{fig:exchange-and-spin}
\end{figure}
